\documentclass[a4paper,11pt]{article}
\usepackage[utf8]{inputenc}
\usepackage[T1]{fontenc}
\usepackage[french]{babel}
\usepackage[left=2cm,right=2cm,top=2cm,bottom=2cm]{geometry}
\usepackage{xcolor}
\usepackage{hyperref}
\usepackage{fontawesome5}
\usepackage{parskip}

\definecolor{primary}{RGB}{0, 50, 130}
\hypersetup{colorlinks=true, urlcolor=primary}

\begin{document}

\noindent
\begin{minipage}[t]{0.5\textwidth}
    \textbf{Lo\"ic Aron MBASSI EWOLO}\\
    \'El\`eve-Ing\'enieur Bac+5 -- Double dipl\^ome ENSEA\\[3pt]
    \faMapMarker\ Cergy, France\\
    \faPhone\ (+33) 7 59 52 33 98\\
    \faEnvelope\ \href{mailto:loicaron.mbassiewolo@ensea.fr}{loicaron.mbassiewolo@ensea.fr}
\end{minipage}
\hfill
\begin{minipage}[t]{0.4\textwidth}
    \raggedleft
    Cergy, le 21 mars 2026
\end{minipage}

\vspace{1.2cm}

\hfill
\begin{minipage}[t]{0.5\textwidth}
    \raggedleft
    \textbf{Centre National d'Etudes Spatiales (CNES)}\\
    Service Recrutement\\
    Toulouse, France
\end{minipage}

\vspace{1cm}

\noindent\textbf{Objet :} Candidature en alternance -- Ing\'enieur IA / ML pour la Surveillance des Satellites -- Septembre 2026

\vspace{0.8cm}

Madame, Monsieur,

La t\'el\'em\'etrie satellitaire g\'en\`ere des flux de s\'eries temporelles multivari\'ees dont l'analyse automatis\'ee par IA -- d\'etection d'anomalies, pr\'ediction de d\'efaillances, classification d'\'ev\'enements de bord -- am\'eliore significativement la s\^uret\'e des missions. Le CNES, agence spatiale nationale, est l'acteur id\'eal pour appliquer l'IA \`a ces probl\'ematiques \`a fort enjeu, o\`u chaque mod\`ele doit \^etre fiable, interpr\'etable et robuste.

Mon exp\'erience sur les s\'eries temporelles biologiques (UROP 2024, collaboration Google DeepMind) est directement analogue : entra\^inement de CNN pour la d\'etection d'arythmies sur signaux ECG multivari\'es, avec optimisation du pipeline d'entra\^inement et m\'etriques cliniques de sensibilit\'e/sp\'ecificit\'e. Mon projet Drone FDI/FTC (MATLAB/Simulink) m'a form\'e \`a la d\'etection de d\'efauts et \`a la tol\'erance aux pannes dans des syst\`emes de contr\^ole -- compl\'ementaire \`a la surveillance de sant\'e satellite. Mon stage \`a l'INRIA Saclay (avril-ao\^ut 2026) apporte la rigueur de la recherche scientifique : pipeline ML valid\'e, m\'etriques de performance documentes, d\'eploiement reproductible.

La formation ENSEA en syst\`emes embarqu\'es et signal, coupl\'ee \`a ma formation en math\'ematiques appliqu\'ees (Polytechnique Yaound\'e), me donne une vision syst\`eme compl\`ete : des contraintes mat\'erielles des capteurs embarqu\'es jusqu'aux algorithmes ML d'analyse. Le CNES pr\'ecise que la ma\^itrise des syst\`emes embarqu\'es temps r\'eel est un plus -- je dispose de cette exp\'erience.

Je suis disponible d\`es septembre 2026 pour cette alternance de 12 mois \`a Toulouse. Appliquer des m\'ethodes ML \`a la surveillance de satellites au CNES est une perspective scientifiquement stimulante et strat\'egiquement importante. Je serais heureux d'en discuter lors d'un entretien.

Je vous prie d'agr\'eer, Madame, Monsieur, l'expression de mes salutations distingu\'ees.

\vspace{1cm}

\noindent\textbf{Lo\"ic Aron MBASSI EWOLO}

\end{document}
